\chapter{Apprentissage en ligne et dérive}\label{ch:ml:online-learning-and-drift}

Un modèle à court horizon est réentraîné chaque week-end. Un mardi, la relation qu'il a apprise change (une plateforme modifie sa grille
tarifaire, un gros intervenant change sa façon de trader, un régime prend fin) et le modèle continue de trader l'ancienne relation jusqu'au
prochain week-end. Un modèle qui se mettrait à jour toutes les heures l'aurait remarqué le mardi après-midi~; un modèle qui se
mettrait à jour à chaque transaction aurait passé toute l'année à courir après le bruit. Ce chapitre traite du choix entre ces deux
échecs~: à quelle vitesse un modèle doit oublier, comment détecter que le monde a changé et quand le réentraîner. Sur un
flux synthétique dont les coefficients changent à des instants aléatoires, un détecteur calibré pour une fausse alarme par an repère un
renversement du signal environ vingt pas après qu'il se produit, et la meilleure mémoire d'un estimateur récursif passe de 200
observations à la totalité d'entre elles à mesure que les régimes s'allongent, de cent pas à dix mille.

\section{Types de dérive}

\begin{definition}[Dérive conceptuelle, décalage des covariables]\label{def:ml:online-learning-and-drift:drift}
La \emph{dérive conceptuelle}\index{dérive conceptuelle} (concept drift) est un changement au cours du temps de la loi conditionnelle de la cible sachant les
variables, $\P_t(y\mid x)$~: la relation que le modèle a apprise cesse de tenir. Le \emph{décalage des covariables}\index{décalage des covariables} (covariate shift)
est un changement de la loi des variables, $\P_t(x)$, avec $\P(y\mid x)$ inchangée.
\end{definition}

Le \omterm{def:ml:online-learning-and-drift:drift}{décalage des covariables} nuit à un modèle qui a appris $\P(y\mid x)$ seulement là où les données d'entraînement étaient denses et doit extrapoler
ailleurs~; un modèle correctement spécifié n'en souffre pas. La \omterm{def:ml:online-learning-and-drift:drift}{dérive conceptuelle} nuit à tous les modèles, et c'est ce que produisent les marchés~:
l'encombrement (crowding) affaiblit un signal (livre~7, chapitre~28), une rupture structurelle l'inverse, un changement de structure de marché redéfinit le
lien entre flux d'ordres et prix. Le livre~7 mesurait après coup la décroissance et les ruptures (tests CUSUM, demi-vie
du signal)~; ce chapitre s'y adapte au fil de l'eau.

Le flux du chapitre compte 20\,000 pas et cinq variables~; la cible est $y_t = \beta_t^\top x_t + \varepsilon_t$ avec un
signal valant 5~\% de la variance, et les coefficients sont tirés à nouveau à la fin de chaque régime, dont l'espacement suit une loi exponentielle de
moyenne $D$ pas. Le R-carré propre à la vérité, le meilleur qu'un modèle puisse atteindre, se situe entre 4.9 et 5.4~\%.

\section{Estimateurs récursifs et oubli}

\begin{definition}[Apprentissage en ligne, moindres carrés récursifs, facteur d'oubli]\label{def:ml:online-learning-and-drift:rls}
L'\emph{apprentissage en ligne}\index{apprentissage en ligne} (online learning) met à jour un modèle à chaque nouvelle observation, à temps et mémoire constants,
et s'évalue de façon préquentielle~: chaque observation est prédite avant d'être utilisée pour apprendre. Les \emph{moindres carrés
récursifs}\index{moindres carrés récursifs} (recursive least squares, RLS) mettent à jour la solution des moindres carrés pondérés
$\hat\beta_t = \arg\min_b\sum_{s\le t}\lambda^{t-s}(y_s - x_s^\top b)^2$ par
\[
k_t = \frac{P_{t-1}x_t}{\lambda + x_t^\top P_{t-1}x_t},\quad \hat\beta_t = \hat\beta_{t-1} + k_t(y_t - x_t^\top\hat\beta_{t-1}),\quad
P_t = \frac{1}{\lambda}\bigl(P_{t-1} - k_tx_t^\top P_{t-1}\bigr),
\]
où le \emph{facteur d'oubli}\index{facteur d'oubli} (forgetting factor) $\lambda\in(0,1]$ atténue le poids des anciennes observations~; leurs poids totalisent
environ $1/(1-\lambda)$, qui est la mémoire effective de l'estimateur.
\end{definition}

\begin{proposition}[L'oubli est un filtre de Kalman à marche aléatoire]\label{prop:ml:online:kalman}
Les \omterm{def:ml:online-learning-and-drift:rls}{moindres carrés récursifs} à \omterm{def:ml:online-learning-and-drift:rls}{facteur d'oubli} $\lambda$ sont le filtre de Kalman (livre~4, chapitre~19) pour le modèle espace-état $y_t =
x_t^\top\beta_t + \varepsilon_t$, $\beta_t = \beta_{t-1} + w_t$, dans lequel l'étape de prédiction gonfle la covariance de l'état
d'un facteur $1/\lambda$ au lieu d'ajouter la covariance de $w_t$~: $P_{t\mid t-1} = P_{t-1\mid t-1}/\lambda$, avec un bruit d'observation
unitaire (Sayed et Kailath, 1994).
\end{proposition}

\begin{proof}
Avec $P_{t\mid t-1} = P_{t-1}/\lambda$ et une variance de bruit unitaire, le gain de Kalman vaut $P_{t\mid t-1}x_t/(1 + x_t^\top P_{t\mid
t-1}x_t) = P_{t-1}x_t/(\lambda + x_t^\top P_{t-1}x_t) = k_t$~; la mise à jour de l'état est celle des \omterm{def:ml:online-learning-and-drift:rls}{moindres carrés récursifs}~; la covariance filtrée
$(I - k_tx_t^\top)P_{t\mid t-1} = (P_{t-1} - k_tx_t^\top P_{t-1})/\lambda$ est le $P_t$ des \omterm{def:ml:online-learning-and-drift:rls}{moindres carrés récursifs}.
\end{proof}

La \Cref{fig:ml:online:forgetting} mesure
le compromis~: les mémoires courtes suivent les régimes mais estiment cinq coefficients à partir de peu de données~; les longues estiment bien
les coefficients des régimes déjà terminés. Avec des régimes de 100 ou 500 pas en moyenne, la meilleure mémoire vaut environ 200
observations ($\lambda = 0.995$), et elle récupère 0.75~\% et 2.44~\% de R-carré sur un maximum de 4.9~\%~; avec des régimes de
2\,000 pas, elle vaut 500 observations (3.71~\%), et avec des régimes de 10\,000 pas il ne faut rien oublier ($\lambda = 1$,
5.21~\% sur 5.29~\%). La descente de gradient stochastique en ligne, à \omterm{def:ml:trees-and-boosting:boosting}{taux d'apprentissage} fixe, est une règle d'oubli rudimentaire à elle seule
(0.57~\%, 2.05~\%, 3.69~\% et 4.84~\%).

\begin{omfigure}
\begin{tikzpicture}
\begin{semilogxaxis}[width=11cm,height=6cm,xlabel={mémoire effective $1/(1-\lambda)$ (observations~; à droite~: $\lambda = 1$)},
  ylabel={$R^2$ préquentiel (\%)},tick label style={font=\small},label style={font=\small},xmin=40,xmax=25000,
  ymin=-2,ymax=6,xtick={50,100,200,500,1000,2000,20000},xticklabels={50,100,200,500,1\,000,2\,000,tout},grid=major,
  grid style={gray!25},legend columns=4,legend style={font=\footnotesize,at={(0.5,-0.25)},anchor=north,draw=none,
  fill=none,/tikz/every even column/.append style={column sep=6pt}},table/col sep=comma]
\addplot[omThm,very thick,mark=*] table[x=memory,y=D100]{figdata/ml/12-online-learning-and-drift/forgetting.csv};
\addplot[omProp,very thick,mark=square*] table[x=memory,y=D500]{figdata/ml/12-online-learning-and-drift/forgetting.csv};
\addplot[omDef,very thick,mark=triangle*] table[x=memory,y=D2000]{figdata/ml/12-online-learning-and-drift/forgetting.csv};
\addplot[black!70,very thick,mark=diamond*] table[x=memory,y=D10000]{figdata/ml/12-online-learning-and-drift/forgetting.csv};
\legend{$D = 100$,$D = 500$,$D = 2\,000$,$D = 10\,000$}
\end{semilogxaxis}
\end{tikzpicture}

\omcaption{R-carré préquentiel des \omterm{def:ml:online-learning-and-drift:rls}{moindres carrés récursifs} (après 1\,000 pas de préchauffage) en fonction de leur mémoire effective, pour
flux dont les coefficients changent en moyenne toutes les $D$ pas. Données~: \texttt{ml\_online.forgetting}.}\label{fig:ml:online:forgetting}
\end{omfigure}

\section{Détection de la dérive}

\begin{definition}[Détecteur de dérive, test de Page--Hinkley, fenêtrage adaptatif]\label{def:ml:online-learning-and-drift:detector}
Un \emph{détecteur de dérive}\index{détecteur de dérive} (drift detector) surveille une statistique d'un modèle en production (ses erreurs, ses gains, la
loi d'une variable) et déclenche une alarme lorsque la loi de cette statistique semble avoir changé. Le
\emph{test de Page--Hinkley}\index{test de Page--Hinkley} (Page--Hinkley test) accumule les écarts de la valeur suivie sous sa moyenne courante,
diminués d'une tolérance $\delta$, et déclenche une alarme lorsque la somme dépasse de plus de $h$ son minimum courant (Page, 1954~; Hinkley, 1971).
Le \emph{fenêtrage adaptatif}\index{fenêtrage adaptatif} (adaptive windowing, ADWIN) conserve une fenêtre de valeurs récentes et en supprime la partie
ancienne dès que une coupure de la fenêtre en deux sous-fenêtres présente des moyennes plus éloignées qu'une borne de type Hoeffding au niveau de confiance $\delta$
(Bifet et Gavaldà, 2007).
\end{definition}

Le test CUSUM du livre~7 (chapitre~13) est le troisième détecteur. Les trois surveillent la même statistique~: le gain
$g_t = \hat y_t\,y_t$ du modèle déployé (sa prévision multipliée par le résultat), standardisé sur la période de préchauffage, dont la moyenne est positive tant que le
modèle a raison et baisse lorsqu'il a tort. Leurs seuils sont fixés, sur un long flux sans changement, à une fausse alarme
par an (une pour 252 pas)~: $h = 9$ pour le CUSUM et le Page--Hinkley (tolérance 0.1), $\delta = 0.0032$ pour l'ADWIN. Puis, sur vingt
flux dont les coefficients changent de signe au pas 2\,000, les détecteurs déclenchent leur première vraie alarme après 20.8 pas en
moyenne pour le CUSUM (médiane 17), 27.4 pour le Page--Hinkley (médiane 18.5) et 42.5 pour l'ADWIN (médiane 26.5)~; sur les 1\,000 pas
précédant l'inversion, ils ont déclenché 74, 78 et 91 fausses alarmes sur l'ensemble des vingt flux, proche des 79
qu'implique leur calibration (\cref{fig:ml:online:alarms}).

\begin{omfigure}
\begin{tikzpicture}
\begin{axis}[width=12cm,height=5.6cm,xlabel={pas},ylabel={moyenne sur 50 pas du gain},tick label style={font=\small},
  label style={font=\small},xmin=1600,xmax=2400,ymin=-1.5,ymax=1.9,grid=major,grid style={gray!25},
  xticklabel style={/pgf/number format/1000 sep={}},legend columns=4,legend style={font=\footnotesize,at={(0.5,-0.3)},
  anchor=north,draw=none,fill=none,/tikz/every even column/.append style={column sep=6pt}},table/col sep=comma]
\addplot[gray,thick] table[x=step,y=running]{figdata/ml/12-online-learning-and-drift/gains.csv};
\addplot[only marks,mark=triangle*,mark size=3pt,omThm,restrict expr to domain={\thisrow{detector}}{0:0}]
  table[x=step,y=level]{figdata/ml/12-online-learning-and-drift/alarms.csv};
\addplot[only marks,mark=square*,mark size=2.6pt,omProp,restrict expr to domain={\thisrow{detector}}{1:1}]
  table[x=step,y=level]{figdata/ml/12-online-learning-and-drift/alarms.csv};
\addplot[only marks,mark=*,mark size=2.6pt,omDef,restrict expr to domain={\thisrow{detector}}{2:2}]
  table[x=step,y=level]{figdata/ml/12-online-learning-and-drift/alarms.csv};
\draw[black!70,thick,dashed] (axis cs:2000,-1.5) -- (axis cs:2000,1.9);
\legend{gain,alarme CUSUM,alarme Page--Hinkley,alarme ADWIN}
\end{axis}
\end{tikzpicture}

\omcaption{Un flux unique dont les coefficients s'inversent à l'étape 2\,000 (en pointillés)~: la moyenne sur 50 étapes du
gain standardisé du modèle déployé, et les alarmes de chaque détecteur, depuis 300 étapes avant l'inversion jusqu'à sa première alarme après celle-ci (marqueurs à
des hauteurs arbitraires, une ligne par détecteur). Données~: \texttt{ml\_online}.}\label{fig:ml:online:alarms}
\end{omfigure}

\section{Calendriers de réentraînement}

\begin{definition}[Calendrier de réentraînement]\label{def:ml:online-learning-and-drift:schedule}
Un \emph{calendrier de réentraînement}\index{calendrier de réentraînement} est la règle qui décide quand un modèle est recalé et sur quelles
données~: selon un calendrier (chaque semaine, sur une fenêtre glissante), sur l'alarme d'un détecteur, en continu (un estimateur en ligne), ou par
une combinaison~; il fait partie du modèle et se valide avec lui.
\end{definition}

\begin{table}[htbp]
\centering\small
\begin{tabular}{@{}>{\raggedright\arraybackslash}p{9.6cm}>{\centering\arraybackslash}p{2.6cm}>{\centering\arraybackslash}p{2.4cm}@{}}
\toprule
calendrier (moindres carrés sur les 500 dernières observations sauf indication contraire) & $R^2$ prequential & recalages\\
\midrule
jamais (ajusté sur les 1\,000 premières étapes) & $-7.81~\%$ & 0\\
toutes les 250 étapes & 2.62~\% & 75\\
sur une alarme de Page--Hinkley (recalé sur les données depuis l'alarme) & $-2.67~\%$ & 2\\
\omterm{def:ml:online-learning-and-drift:rls}{moindres carrés récursifs}, $\lambda = 0.998$ & 3.24~\% & à chaque étape\\
\midrule
la vérité & 4.62~\% & \\
\bottomrule
\end{tabular}
\caption{Quatre \omterm{def:ml:online-learning-and-drift:schedule}{calendriers de réentraînement} sur un flux dont les coefficients sont redessinés tous les 2\,000 étapes en moyenne (19\,000 étapes
évaluées après une période de chauffe de 1\,000). Données~: \texttt{ml\_online.schedules}.}\label{tab:ml:online:schedules}
\end{table}

Le \Cref{tab:ml:online:schedules} livre une leçon sur les détecteurs. Le \omterm{def:ml:online-learning-and-drift:detector}{test de Page--Hinkley}, qui a repéré un renversement de signe
en 18.5 étapes médianes, a recalé le modèle deux fois alors que les coefficients ont changé dix fois~: quand les coefficients sont redessinés plutôt que d'être inversés, les
gains de l'ancien modèle tombent à zéro plutôt qu'en dessous, un changement plus faible que celui qu'il était réglé pour redouter, et il attend. Le calendrier
fixe, simple et régulier, le bat~; les \omterm{def:ml:online-learning-and-drift:rls}{moindres carrés récursifs}, avec la mémoire qui convient aux régimes, les battent tous deux. Un détecteur
doit être réglé sur le changement qui compte, et un calendrier qui s'appuie sur lui a besoin d'un calendrier de fond sous-jacent.

\begin{method}[Choisir comment un modèle s'adapte]\label{met:ml:online:method}
\begin{enumerate}
\item Estimer la vitesse à laquelle la relation évolue~: ajuster le modèle sur des fenêtres glissantes de plusieurs longueurs et comparer leurs
scores prequential, ou, de manière équivalente, balayer le \omterm{def:ml:online-learning-and-drift:rls}{facteur d'oubli}.
\item Évaluer prequentiellement chaque calendrier sur l'historique, en tenant compte des coûts de recalage (validation, déploiement, rotation).
\item Calibrer les détecteurs sur des périodes sans changement, pour obtenir un taux de fausses alarmes que le desk peut exploiter, et mesurer
leur délai sur les changements qui comptent (un renversement, un retour à zéro, un changement de volatilité).
\item Combiner~: un recalage calendaire comme plancher, un détecteur pour recaler plus tôt, et un estimateur en ligne là où le modèle est
assez simple pour être mis à jour en toute sécurité.
\end{enumerate}
\end{method}

\section{Tutoriel~: la demi-vie d'un modèle}\label{tut:ml:online-learning-and-drift:tut}

\begin{tutorial}
\textbf{Objectif.} Balayer le \omterm{def:ml:online-learning-and-drift:rls}{facteur d'oubli} des \omterm{def:ml:online-learning-and-drift:rls}{moindres carrés récursifs} en fonction de la longueur des régimes, calibrer trois détecteurs
pour une fausse alarme par an et mesurer leur délai sur un renversement, puis comparer les \omterm{def:ml:online-learning-and-drift:schedule}{calendriers de réentraînement}. \textbf{État final~:}
\cref{fig:ml:online:forgetting,fig:ml:online:alarms}, \cref{tab:ml:online:schedules}.
\begin{tutsteps}
\item \textbf{\omterm{def:ml:online-learning-and-drift:rls}{Moindres carrés récursifs} avec oubli} (\cref{prop:ml:online:kalman}).
\omcode{code/firm/onlinelearn/firm_onlinelearn.py}{51}{67}{Moindres carrés récursifs avec facteur d'oubli.}\label{lst:ml:online:rls}
\item \textbf{Le \omterm{def:ml:online-learning-and-drift:detector}{test de Page--Hinkley}.}
\omcode{code/firm/onlinelearn/firm_onlinelearn.py}{104}{123}{Le test de Page--Hinkley pour une baisse de la moyenne.}\label{lst:ml:online:ph}
\item \textbf{Un ADWIN à fenêtre.}
\omcode{code/firm/onlinelearn/firm_onlinelearn.py}{126}{155}{Fenêtrage adaptatif, vérifié toutes les quelques étapes.}\label{lst:ml:online:adwin}
\item \textbf{Exécuter} \texttt{ml\_online.forgetting(D)}, \texttt{thresholds()}, \texttt{delays()}, \texttt{schedules()} et
\texttt{fig\_online.py}.
\end{tutsteps}
\textbf{Pour aller plus loin.} Faire en sorte que les régimes ramènent leurs coefficients à zéro au lieu de les redessiner, et recalibrer les
détecteurs pour ce changement~; combiner le calendrier avec le détecteur.
\end{tutorial}

\section{Construction~: apprentissage en ligne et détection de dérive}\label{bld:ml:online-learning-and-drift:onlinelearn}

\begin{build}
\bfield{Rôle} Des modèles qui s'adaptent à une vitesse choisie, des détecteurs dont le taux de fausses alarmes est connu, et des calendriers validés comme
les modèles qu'ils recalent.
\bfield{Interface} \texttt{drift\_stream(n, p, mean\_regime, r2, seed, flip\_at)}~; \texttt{RLS(p, lam, delta)},
\texttt{OnlineSGD(p, lr)}, \texttt{prequential(model, X, y)}~; \texttt{CUSUM(mu0, k, h)}, \texttt{PageHinkley(delta, lam)},
\texttt{ADWIN(delta, max\_window, check, min\_side)}, \texttt{alarm\_times},
\texttt{calibrate(make\_for, thresholds, null\_values, rate)}~; \texttt{retrain\_schedule(X, y, policy, every, window, warm, make\_detector, min\_fit)}.
\bfield{Règles} Prédire avant d'apprendre~; les détecteurs ne voient que les valeurs passées et présentes~; les seuils proviennent de la calibration sur
des données sans changement, jamais de la période qu'on évalue.
\bfield{Tests d'acceptation} \texttt{code/firm/onlinelearn/tests/}~: les \omterm{def:ml:online-learning-and-drift:rls}{moindres carrés récursifs} avec $\lambda = 1$ égalent les moindres carrés ordinaires par lots~; leur mémoire effective
correspond à $1/(1-\lambda)$ sur un changement planté~; le taux de fausses alarmes de chaque détecteur sur des données sans changement reste dans
sa calibration, et il repère un décalage planté de la moyenne~; le calendrier fixe recale aux bonnes étapes.
\bfield{Pour aller plus loin} Un oubli qui s'adapte à la statistique du détecteur (variable $\lambda$)~; des détecteurs sur les
distributions des variables (chapitre~27).
\end{build}

\begin{omsources}
\item E.~S.~Page, ``Continuous inspection schemes'', \emph{Biometrika} 41(1/2), 1954.
\item A.~Bifet et R.~Gavaldà, ``Learning from time-changing data with adaptive windowing'', \emph{SIAM International
Conference on Data Mining}, 2007.
\item J.~Gama, I.~Žliobaitė, A.~Bifet, M.~Pechenizkiy et al., ``A survey on concept drift adaptation'', \emph{ACM
Computing Surveys} 46(4), 2014.
\item A.~H.~Sayed et T.~Kailath, ``A state-space approach to adaptive RLS filtering'', \emph{IEEE Signal Processing
Magazine} 11(3), 1994.
\item D.~V.~Hinkley, ``Inference about the change-point from cumulative sum tests'', \emph{Biometrika} 58(3), 1971.
\item G.~Widmer et M.~Kubat, ``Learning in the presence of concept drift and hidden contexts'', \emph{Machine Learning} 23,
1996.
\end{omsources}

\section{Exercices}

\begin{exercise}[$\star$]\label{exo:ml:online-learning-and-drift:1}
Calculer les mémoires effectives des \omterm{def:ml:online-learning-and-drift:rls}{facteurs d'oubli} 0.99, 0.998 et 0.9995. Calculer le poids
qu'une observation vieille de 500 étapes reçoit avec chacun d'eux, par rapport à la plus récente.
\end{exercise}

\begin{exercise}[$\star$]\label{exo:ml:online-learning-and-drift:2}
Un détecteur lève une fausse alarme toutes les 252 étapes. Calculer le nombre de fausses alarmes que doivent lever 20 flux de 1\,000 étapes, puis comparer
avec les décomptes du chapitre.
\end{exercise}

\begin{exercise}[$\star$]\label{exo:ml:online-learning-and-drift:3}
Un CUSUM de référence 0 et de tolérance $k = 0.1$ suit des gains standardisés dont la moyenne passe de 0 à $-0.38$ au
retournement (les flux du chapitre). À quel rythme moyen sa statistique croît-elle après le changement, combien d'étapes lui faut-il depuis zéro
pour franchir $h = 9$, et pourquoi le délai médian mesuré est-il plus court~?
\end{exercise}

\begin{exercise}[$\star\star$]\label{exo:ml:online-learning-and-drift:4}
Pourquoi la meilleure mémoire croît-elle avec la longueur du régime, et pourquoi n'est-elle pas proportionnelle à celle-ci~?
\end{exercise}

\begin{exercise}[$\star\star$]\label{exo:ml:online-learning-and-drift:5}
Pourquoi le calendrier de Page--Hinkley n'a-t-il réentraîné le modèle que deux fois, et que changeriez-vous~?
\end{exercise}

\begin{exercise}[$\star\star$]\label{exo:ml:online-learning-and-drift:6}
\emph{Trouver la faille.} «~Nous avons choisi le seuil du détecteur pour qu'il capte les deux changements de régime de l'an dernier en une semaine,
et il l'aurait fait.~»
\end{exercise}

\begin{exercise}[$\star\star\star$]\label{exo:ml:online-learning-and-drift:7}
\emph{Programmation.} Exécuter \texttt{ml\_online.schedules} avec des régimes de 500 étapes en moyenne (\texttt{schedules(D=500)}). Présenter les
quatre calendriers et expliquer ce qui a changé.
\end{exercise}

\begin{exercise}[$\star\star\star$]\label{exo:ml:online-learning-and-drift:8}
Montrer que les \omterm{def:ml:online-learning-and-drift:rls}{moindres carrés récursifs} (RLS) avec $\lambda = 1$ et $P_0 = \delta I$ donnent l'estimateur ridge de pénalité $1/\delta$, et donc les moindres carrés ordinaires
lorsque $\delta\to\infty$.
\end{exercise}

\section{Problème~: la demi-vie d'un modèle}

\begin{problem}[{Problème du week-end --- à quelle vitesse oublier}]\label{pb:ml:online-learning-and-drift:1}
Le flux du chapitre~: cinq variables, un signal valant 5~\% de la variance, des coefficients tirés à nouveau à chaque $D$ étapes en moyenne.

\medskip\noindent\textbf{Partie I --- Oubli.}
\begin{enumerate}
\item Quel est le R-carré de la vérité, et pourquoi aucun modèle n'en est-il proche pour les régimes courts~?
\item Quel \omterm{def:ml:online-learning-and-drift:rls}{facteur d'oubli} est le meilleur pour $D = 100$, 500, 2\,000 et 10\,000, et quel score obtient-il~?
\item Quel score obtient la descente de gradient stochastique en ligne, et pourquoi se comporte-t-elle comme un oubli~?
\item Que dit la \cref{prop:ml:online:kalman} de l'hypothèse qu'un \omterm{def:ml:online-learning-and-drift:rls}{facteur d'oubli} suppose implicitement~?
\end{enumerate}

\medskip\noindent\textbf{Partie II --- Détecteurs.}
\begin{enumerate}[resume]
\item Quelle statistique les détecteurs surveillent-ils, et pourquoi le gain plutôt que l'erreur quadratique~?
\item Quels seuils donnent une fausse alarme par an~?
\item Quels sont leurs délais moyen et médian après un retournement~?
\item Combien de fausses alarmes ont-ils levées avant le retournement, et est-ce trop~?
\end{enumerate}

\medskip\noindent\textbf{Partie III --- Calendriers.}
\begin{enumerate}[resume]
\item Que donnent les quatre calendriers avec des régimes de 2\,000 étapes~?
\item Pourquoi le modèle jamais réentraîné obtient-il un score négatif~?
\item Pourquoi le calendrier piloté par le détecteur a-t-il échoué~?
\item Pourquoi les \omterm{def:ml:online-learning-and-drift:rls}{moindres carrés récursifs} l'emportent-ils~?
\end{enumerate}

\medskip\noindent\textbf{Partie IV --- Verdict.}
\begin{enumerate}[resume]
\item Énoncer le \emph{résultat nommé}~: le \omterm{def:ml:online-learning-and-drift:rls}{facteur d'oubli} qui minimise la perte hors échantillon en fonction de la
longueur du régime, et le délai de chaque détecteur à une fausse alarme par an.
\item Quel calendrier retenir pour un modèle à horizon court sur une plateforme changeante~?
\item Comment estimer $D$ sur des données réelles~?
\item Que coûte une fausse alarme, et que coûte un changement manqué~?
\item En quoi ce chapitre se rattache-t-il à la demi-vie du signal du livre~7~?
\item Que doit consigner un système de suivi pour chaque alarme~?
\item Comment le \omterm{def:ml:online-learning-and-drift:drift}{décalage des covariables} apparaîtrait-il dans ces statistiques~?
\item En une phrase~: à quelle vitesse un modèle doit-il oublier~?
\end{enumerate}
\end{problem}

\section{Questions d'entretien}

\begin{interviewq}[$\star$ \iqroles{ingénieur ML, chercheur}]\label{iq:ml:online-learning-and-drift:1}
Quelle est la différence entre \omterm{def:ml:online-learning-and-drift:drift}{dérive conceptuelle} et \omterm{def:ml:online-learning-and-drift:drift}{décalage des covariables}~? Donnez un exemple de marché pour chacune.
\end{interviewq}

\begin{interviewq}[$\star\star$ \iqroles{chercheur}]\label{iq:ml:online-learning-and-drift:2}
Dérivez les \omterm{def:ml:online-learning-and-drift:rls}{moindres carrés récursifs} avec un \omterm{def:ml:online-learning-and-drift:rls}{facteur d'oubli}, et dites ce que signifie ce facteur.
\end{interviewq}

\begin{interviewq}[$\star\star$ \iqroles{ingénieur ML}]\label{iq:ml:online-learning-and-drift:3}
Comment fixez-vous le seuil d'un \omterm{def:ml:online-learning-and-drift:detector}{détecteur de dérive}, et comment savez-vous qu'il fonctionne~?
\end{interviewq}

\begin{interviewq}[$\star\star$ \iqroles{chercheur, trader}]\label{iq:ml:online-learning-and-drift:4}
Votre modèle est réentraîné chaque mois. Comment décideriez-vous s'il faut le réentraîner chaque semaine, chaque jour, ou en continu~?
\end{interviewq}

\begin{interviewq}[$\star\star$ \iqroles{ingénieur ML}]\label{iq:ml:online-learning-and-drift:5}
Expliquez le \omterm{def:ml:online-learning-and-drift:detector}{test de Page--Hinkley} et l'une de ses faiblesses.
\end{interviewq}

\begin{interviewq}[$\star\star\star$ \iqroles{chercheur}]\label{iq:ml:online-learning-and-drift:6}
Montrez que l'oubli exponentiel est un filtre de Kalman pour un coefficient en marche aléatoire, et ce que cela implique pour le choix de
$\lambda$.
\end{interviewq}
